\documentclass[conference]{IEEEtran}
\IEEEoverridecommandlockouts

\usepackage{cite}
\usepackage{amsmath,amssymb,amsfonts}
\usepackage{algorithm}
\usepackage{algpseudocode}
\usepackage{graphicx}
\usepackage{textcomp}
\usepackage{xcolor}
\usepackage{booktabs}
\usepackage{url}
\usepackage[hidelinks]{hyperref}

\def\BibTeX{{\rm B\kern-.05em{\sc i\kern-.025em b}\kern-.08em T\kern-.1667em\lower.7ex\hbox{E}\kern-.125emX}}

\begin{document}

\title{Coordination Before Trajectories: Selecting Multi-Robot Manoeuvres\\
from the Geometry of the Conflict Graph}

% TODO: fill in author block before submission. Square brackets must stay braced
% here -- a bare [ after \\ is parsed as the optional length argument of \\.
\author{\IEEEauthorblockN{{Author Name}}
\IEEEauthorblockA{\textit{{Department}} \\
\textit{{Institution}} \\
{City, Country} \\
{email@institution.edu}}
}

\maketitle

\begin{abstract}
Kinodynamic multi-robot motion planners overwhelmingly resolve inter-robot conflicts by
\emph{searching harder}: a conflict is detected, a constraint is imposed, and some
planner---a trajectory optimiser, a sampling-based planner, or a joint-space solver---is
asked again for a trajectory. As robot density rises, the number of conflicts rises with
it, and the planner spends its budget re-answering the same question. We argue that at
density the cheaper question is not \emph{``find me a trajectory''} but \emph{``what is the
coordination?''}, and that the answer is already latent in the geometry of the conflict
graph. We present a constructive planner that reads a coordination \emph{device} off each
connected component of the conflict graph---lanes when the component is a set of distinct
crossings spread through space, a roundabout when every route in the component passes
through a single point---assigns each robot a lane or an orbit, realises the result as
rest-to-rest legs under the true acceleration bounds, schedules departures by insertion,
and verifies the outcome with the simulator's own collision checker before committing.
On four of five dense benchmark scenarios with 16--32 second-order unicycles the planner
returns a verified, collision-free plan that reaches every goal with \textbf{zero solver
calls}, in 7--51\,s of wall time. The fifth is analysed rather than claimed: we show its
conflict component demands more lanes than its corridor geometry can hold, and we retract
two earlier diagnoses of it that measurement did not support. We are explicit about what
the construction gives up: conflicts are detected on a discretised progress
parameterisation, and the realised trajectories fire linear and angular acceleration in
separate phases, costing 1.70--2.08$\times$ the cruise-limited lower bound.
\end{abstract}

\begin{IEEEkeywords}
multi-robot systems, kinodynamic motion planning, coordination, conflict resolution
\end{IEEEkeywords}

\section{Introduction}

Teams of mobile robots increasingly share tight physical space: order-picking robots in
warehouse aisles, container movers on quaysides, and inspection platforms in greenhouse
rows all operate in environments where the free space is a small multiple of the robot
body. In such settings the interesting difficulty is rarely finding a path for any one
robot---a single robot's route is often almost a straight line---but deciding how the
robots are to get past \emph{one another}.

Two properties make this hard. The first is combinatorial: the joint configuration space
of $n$ robots grows exponentially in $n$, so planning in it directly is impractical beyond
a handful of robots. The second is dynamic. Multi-agent path finding (MAPF) in its
classical form \cite{cbs2015} assumes agents occupy grid cells and may wait in place at no
cost; a robot with second-order dynamics can do neither. It cannot stop instantly, cannot
turn without either stopping or sweeping an arc whose radius is set by its speed, and
cannot be assumed to arrive anywhere at a time chosen by the planner. Kinodynamic
multi-robot motion planning (MRMP) is the problem of respecting both.

\subsection{The prevailing paradigm and its cost}

The dominant response to the combinatorial difficulty is to decouple. Each robot is
planned for individually; the resulting trajectories are checked against one another;
where they conflict, a constraint is imposed and some planner is invoked again. This
structure is shared, with variations, across a large literature: prioritized planning
orders the robots and plans each against its predecessors \cite{erdmann1987}; conflict-based
search builds a tree of constraint sets and re-plans single agents at the leaves
\cite{cbs2015}; and the kinodynamic descendants of that idea replace the single-agent
planner with a dynamics-aware one---discontinuity-bounded search over motion primitives in
db-CBS \cite{dbcbs2024}, and a lightweight MAPF layer over the same primitives in
db-LaCAM \cite{dblacam2025}. K-ARC \cite{karc2025} interleaves optimisation-based and
sampling-based planning, constructing its solution iteratively in segments and escalating
to more expensive machinery on the segments that resist.

What every one of these has in common is the \emph{unit of work}. Each conflict, once
detected, is discharged by asking a planner for a trajectory. The planner is good; the
question is expensive; and crucially, the number of times it is asked grows with the number
of conflicts, which grows with density. In our own faithful reimplementation of the K-ARC
ladder, a 32-robot instance of a dense crossing benchmark consumed over 300 solver calls
and the entire 1800\,s budget without returning a valid plan. The failure is not a bug in
the rungs. It is that at this density the planner is answering the wrong question hundreds
of times.

\subsection{An observation about dense scenarios}

Consider what a human traffic engineer does with the same problem. Faced with a two-way
street, they do not compute a trajectory for each vehicle; they paint a line down the
middle and assign a side. Faced with a junction where many roads meet at one point, they
do not resolve the crossings pairwise; they build a roundabout and impose a single
circulation direction. Both are \emph{conventions}, agreed before any vehicle moves, and
both work because they exploit a structural property of the situation---in the first case
that the conflict is a two-sided encounter, in the second that many routes converge on one
place.

The observation motivating this paper is that these structural properties are readable
directly off the conflict graph, and that different components of the same problem may
warrant different conventions. A conflict component whose pairwise encounters are scattered
through space is a set of independent two-sided crossings, and lanes resolve it. A
component whose encounters all pile onto a single point is a junction, and no number of
sides resolves it---there are only two sides, and there may be thirty-two routes. That
component needs a roundabout.

This is not merely an analogy. In the benchmarks we study, the two situations are separated
cleanly by a single measurable quantity, the spatial spread of the component's pairwise
closest-approach points, and the separation is not marginal: it is $4.6$--$4.9$\,m in the
scenarios that want lanes and exactly $0.0$\,m in the scenarios that want a roundabout.

\subsection{Contributions}

\begin{enumerate}
\item \textbf{Device selection from conflict geometry.} We show that a connected component
of the conflict graph carries enough information to choose its own coordination
mechanism, and give a concrete test---the concentration of pairwise closest-approach
points---that separates ``several distinct crossings'' from ``one junction''.

\item \textbf{Two devices, both constructive.} A \emph{lane} device that colours the
conflict graph and displaces routes along a shared world-frame axis, with the lane count
and pitch bounded by the corridor that physically exists; and a \emph{roundabout} device
that displaces every member of a junction to its own right onto a ring of radius
$|C|(d_{\mathrm{diag}}{+}\varepsilon)/2\pi$, producing one shared circulation. Neither
device calls a solver.

\item \textbf{An insertion schedule with two targeted repairs}, motivated by measured
failure modes rather than by generic search: a precedence order derived from where robots
come to rest, and a single-blocker eviction step.

\item \textbf{A verify-or-fall-through contract.} Nothing is committed that does not pass
the simulator's own collision checker, so the construction can only add solutions to an
existing ladder, never remove one.

\item \textbf{An honest failure analysis.} We report the one benchmark the method does not
solve, identify its cause as a lane-capacity shortfall, and retract two earlier
explanations of it that further measurement contradicted---including one produced by a
diagnostic of our own that was subtly mis-specified.
\end{enumerate}

\section{Related Work}

\subsection{Multi-agent path finding}
CBS \cite{cbs2015} is the reference point for conflict-driven multi-agent planning. It
searches a binary constraint tree at the high level and re-plans individual agents at the
low level, and is optimal for the discrete problem. Its assumptions---unit-time actions,
cost-free waiting, agents as cells---are precisely the ones that a dynamically constrained
robot violates, which is what motivates the kinodynamic variants below.

\subsection{Single-robot kinodynamic planning}
Randomised kinodynamic planning \cite{lavalle2001} established sampling in state space
under a dynamics model as the general-purpose answer for a single robot, and remains the
low-level workhorse in many multi-robot systems, including the sampling rungs of the
baseline we compare against.

\subsection{Kinodynamic multi-robot planning}
db-CBS \cite{dbcbs2024} combines CBS with discontinuity-bounded A$^*$ over precomputed
motion primitives, resolving inter-robot collisions by constraining the primitive-level
search and finishing with a joint-space trajectory optimisation. db-LaCAM
\cite{dblacam2025} keeps the primitive representation but replaces the conflict-tree layer
with a lightweight MAPF method, trading optimality for scale. K-ARC \cite{karc2025} plans
in segments, interleaving optimisation and sampling, and reports scalability to 32 robots.
Our baseline is a reimplementation of the K-ARC ladder, since the authors have not released
code.

\subsection{Where this work sits}
All of the above resolve a conflict by re-invoking a planner under new constraints. We are
not aware of prior work that treats the conflict component's \emph{geometry} as a selector
over qualitatively different coordination mechanisms, and derives the manoeuvre
analytically once the mechanism is chosen. The devices themselves are not novel in
isolation---lane discipline and roundabouts are old ideas in traffic engineering---but
choosing between them per conflict component, from a measurable property of that
component, and realising the choice under true second-order bounds, is the contribution.

\section{Problem Formulation}

We consider $n$ robots sharing a bounded planar workspace $\mathcal{W} \subset
\mathbb{R}^2$ containing static obstacles $\mathcal{O}$. Each robot $i$ has state
$\mathbf{x}_i = [x, y, \theta, v, \omega]^\top$ and control
$\mathbf{u}_i = [a, \alpha]^\top$, evolving under the second-order unicycle
\begin{equation}
\dot{\mathbf{x}} =
\begin{bmatrix} v\cos\theta & v\sin\theta & \omega & a & \alpha \end{bmatrix}^\top ,
\label{eq:dyn}
\end{equation}
subject to $|v| \le v_{\max}$, $|\omega| \le \omega_{\max}$, $|a| \le a_{\max}$ and
$|\alpha| \le \alpha_{\max}$. The body is a rectangle of forward extent $w$ and lateral
extent $\ell$; we write $d_{\mathrm{diag}} = \sqrt{w^2 + \ell^2}$ for its diagonal, which
is the centre distance at which two bodies can touch in the worst relative orientation.

Given starts $\mathbf{x}_i(0)$ and goals $\mathbf{g}_i$, we seek controls
$\mathbf{u}_i(\cdot)$ over a horizon $T$ such that every robot reaches its goal region,
$\|\mathbf{p}_i(T) - \mathbf{g}_i\| < r_{\mathrm{goal}}$; no robot intersects an obstacle;
no two bodies come within a clearance $\varepsilon$ of one another at any instant; and
$T$ does not exceed the episode budget $T_{\max}$.

Two remarks fix the standard we hold ourselves to. First, the collision condition is a
\emph{clearance}, not merely non-overlap: we require the surface-to-surface distance
between any two bodies to stay above $\varepsilon$, which is strictly stronger than
requiring them not to intersect. Second, the acceptance test is the simulator's own
checker applied to the executed rollout, not the planner's internal model of it.

\section{Method}

\subsection{Overview}

The planner is constructive: it decides the coordination first and derives the
trajectories from it. Algorithm~\ref{alg:main} gives the structure. There is no search
loop, no constraint tree, and no call to any optimiser or sampling-based planner.

\begin{algorithm}[t]
\caption{Conflict-derived constructive planning}
\label{alg:main}
\begin{algorithmic}[1]
\Require starts, goals, optional kinematic guides
\Ensure a verified plan, or $\varnothing$
\State $R \gets$ \textsc{References}(starts, goals, guides)
\State $P \gets \{(i,j) : \mathrm{sep}(R_i,R_j) < d_{\mathrm{diag}} + \varepsilon\}$
\ForAll{connected components $C$ of $(V,P)$}
  \If{\textsc{IsJunction}($C$)} \Comment{Sec.~\ref{sec:device}}
    \State assign each $u \in C$ a right-hand \textbf{orbit} \Comment{Sec.~\ref{sec:orbit}}
  \Else
    \State assign each $u \in C$ a \textbf{lane} \Comment{Sec.~\ref{sec:lane}}
  \EndIf
\EndFor
\State displace and simplify each route; build rest-to-rest legs
\State $\delta \gets$ \textsc{InsertionSchedule}(tracks) \Comment{Sec.~\ref{sec:sched}}
\State \Return \textsc{Verify}(tracks, $\delta$) $?$ plan $:$ $\varnothing$
\end{algorithmic}
\end{algorithm}

\subsection{Reference routes and progress parameterisation}
\label{sec:refs}

Each robot's intended route is sampled at $m = 64$ points spaced evenly by
\emph{arclength}, giving a normalised progress parameter $s \in [0,1]$. With a kinematic
guide available---any obstacle-free path, in our experiments from a shortcut RRT---the
reference follows it; without one it degrades to the straight start--goal segment.

Arclength sampling is what makes index $k$ comparable across robots whose routes differ in
length: each robot is $k/m$ of the way through its own journey. This is a deliberate and
consequential approximation, and we state its content plainly. Comparing robots at equal
\emph{progress} rather than equal \emph{time} amounts to assuming that they advance along
their routes at comparable rates. Where that assumption holds, the conflict set computed on
this grid is a good proxy for the set of genuine spatio-temporal conflicts; where it fails,
the proxy both misses conflicts and invents them. The reference is not a plan and is never
executed: it exists only to answer \emph{who contends, and where along each route}. The
trajectories themselves, described in Sec.~\ref{sec:legs}, are continuous-time
integrations of \eqref{eq:dyn}. Sec.~\ref{sec:lim} returns to this as a limitation.

\subsection{The conflict graph}

Robots $i$ and $j$ conflict if their references come within
$d_{\mathrm{diag}} + \varepsilon$ at the same progress. The diagonal, rather than the
lateral extent, is the correct threshold: two rectangles can touch with their centres a
diagonal apart, and a pair missed at this stage never receives a coordination assignment at
all. We take connected components of the resulting graph; each is treated independently.

\subsection{Device selection}
\label{sec:device}

For a component $C$ with $|C| \ge 3$, let $H(C)$ be the set of pairwise closest-approach
midpoints of its members' references. We call $C$ a \emph{junction} if all of $H(C)$ lies
within a body-sized ball,
\begin{equation}
\max_{h \in H(C)} \| h - \bar{h} \| \;<\; d_{\mathrm{diag}} + \varepsilon ,
\qquad \bar{h} = \tfrac{1}{|H|}\textstyle\sum_{h \in H} h .
\label{eq:hub}
\end{equation}
Components of size two are never junctions: a lane already separates two robots.

The intuition is that a lane is a \emph{two-sided} device. It resolves a crossing by
putting one robot on each side of it, which is exactly right when a component is a handful
of encounters distributed through space, and structurally incapable of separating $|C|$
routes that all pass through one point. Sec.~\ref{sec:res-device} shows this test is not
delicate in practice: the quantity in \eqref{eq:hub} is $4.6$--$4.9$\,m for our lane
scenarios and exactly $0.0$\,m for our junction scenarios.

\subsection{The lane device}
\label{sec:lane}

\paragraph{Sides} We greedily colour the conflict graph so that conflicting robots receive
different colours; the colour indexes a lane.

\paragraph{Axis} Lanes are displacements along a shared \emph{world-frame} axis, computed
per group from the dominant direction of member tangents by eigendecomposition of
$\sum_k \mathbf{t}_k \mathbf{t}_k^\top$, taking the normal of the leading eigenvector. A
per-robot body-frame normal is wrong here, and instructively so: it flips with heading, so
two head-on robots given opposite lane signs are displaced to the \emph{same} side of the
road.

\paragraph{Direction groups} One axis per connected component is right when the component
\emph{is} one encounter, and wrong when it has grown by transitivity. We therefore split a
component by tangent orientation, taken modulo $\pi$ so that a head-on pair---whose
tangents differ by exactly $\pi$---stays in one group and keeps the shared axis its
opposite lane signs depend on, while genuinely differently-oriented encounters are
separated.

\paragraph{Capacity} The colouring may demand more lanes than the workspace holds. Let
$\rho$ be the closest a robot's reference comes to a reference it does \emph{not} contend
with; the room available for displacement is $\rho - (\ell + \varepsilon)$. We cap the
lane count at $\lfloor \rho'/(\ell{+}\varepsilon)\rfloor + 1$ with
$\rho' = \rho - (\ell + \varepsilon)$, and set the pitch to
$\min(w_{\mathrm{lane}}, \rho'/(L{-}1))$, compressing only when the nominal width does not
fit. Without this cap the outer lanes of a wide colouring reach into the neighbouring
corridor, and the device manufactures conflicts instead of resolving them
(Sec.~\ref{sec:failure}).

\paragraph{Realisation} The displacement is applied over the progress window in which the
robot is contended, padded and tapered at both ends so start and goal are untouched. If the
assigned lane is blocked by an obstacle the robot \emph{changes lane}---every lane position
is tried, nearest to the assigned one first---rather than reverting to its centre line,
where its head-on partner already is. Reverting leaves the encounter to be resolved in
time, which is the expensive resolution.

\subsection{The roundabout device}
\label{sec:orbit}

For a junction component with centre $\bar{h}$, every member is displaced to its
\emph{own right} onto a ring of radius
\begin{equation}
\rho_{\mathrm{ring}} \;=\; \frac{|C| \, (d_{\mathrm{diag}} + \varepsilon)}{2\pi},
\label{eq:ring}
\end{equation}
which is the radius at which $|C|$ bodies would fit around the ring at the required
spacing if they all arrived simultaneously. Concretely, if $\mathbf{t}$ is the robot's unit
tangent at closest approach, the displacement axis is the right-hand normal
$[t_y, -t_x]^\top$ and the magnitude is whatever brings the route to $\rho_{\mathrm{ring}}$
from $\bar{h}$.

The right-hand rule is what makes the displacement \emph{collective}. A shared axis cannot
express it, because ``the same side'' of a junction is a different direction for each
approach: a robot heading $+x$ must pass below the centre while one heading $+y$ passes to
its right, and only a per-robot right-hand normal makes those two senses agree into a single
circulation. This is the exact converse of the lane case, where a shared axis is essential
and a per-robot normal is wrong---which is why the two devices cannot be collapsed into one.

Mirroring is available to a lane as a fallback but never to an orbit: sending one member the
other way round a junction puts it head-on into the entire circulation.

\subsection{Trajectory realisation}
\label{sec:legs}

The displaced route is simplified by Douglas--Peucker \cite{douglas1973} on deviation from
the chord. An angle-threshold simplifier is not usable here: a lane displacement is gradual
by construction, so its per-step turn falls below any threshold and the bulge is deleted
outright.

Each remaining waypoint is realised as a rest-to-rest leg: turn in place to face the
waypoint, then drive to it. Each phase is a one-dimensional trapezoidal profile---accelerate
at the bound, cruise at the velocity bound if the distance warrants it, decelerate at the
bound---and each is integrated through the same dynamics \eqref{eq:dyn} and the same
integrator the simulator uses, so the leg is feasible by construction rather than by
assertion.

This decomposition is what makes closed-form realisation possible: turning in place has
$v = 0$, so position is invariant, and driving straight has $\omega = 0$, so heading is
invariant, and the two one-dimensional profiles do not interact. It also guarantees that the
executed path \emph{is} the polyline on which the lane and clearance reasoning was done. It
costs a factor we measure and report in Sec.~\ref{sec:lim}.

\subsection{Insertion scheduling}
\label{sec:sched}

Lane assignment is spatial and can be decided from routes alone. \emph{When} a robot moves
cannot: two robots share a place only if they are there at the same time, and the duration
of a leg is known only once it is built.

We schedule by insertion. Robots are considered in some order; each is given the least
departure delay at which it clears every robot already placed. The clash test compares the
two trajectories over their full padded span, because a robot occupies its start until it
departs and its goal forever after---omitting the tail permits a plan in which a robot
parks on a spot a later robot drives through.

Insertion is greedy and its only lever is delay, which produces a characteristic failure:
when a robot must be \emph{gone before} another arrives, waiting is exactly the wrong
move, and the run dead-ends with no way back. We add two repairs, each motivated by a
measured failure rather than by generic backtracking.

\paragraph{Rest-derived precedence} Robot $r$ is ordered before $q$ if $q$'s goal lies on
$r$'s route (so $q$ parked would block $r$) or if $r$'s start lies on $q$'s route (so $r$
must clear out before $q$ arrives). Swaps induce two-cycles---each partner's goal is the
other's start---which is not a contradiction, since both simply leave early; cycles are
broken by route length and mutual partners are kept adjacent in the order.

\paragraph{Single-blocker eviction} If a robot cannot be seated, we find the delay at which
it is least obstructed. If exactly one already-placed robot stands in the way there, that
robot is evicted, the stranded robot is seated, and the evicted robot is re-queued. Each
robot may be evicted a bounded number of times.

Both repairs are applied only \emph{after} plain insertion has been tried under every
candidate order, because each one widens the feasible set and thereby changes which
solution the search returns first---and the rescued schedule is not necessarily the better
one. Sec.~\ref{sec:res-sched} quantifies this.

\subsection{Verification and the fallback contract}

The assembled plan is checked with the simulator's own collision routines: every pair at
every shared time index, every robot against every obstacle, every terminal position
against the true goal radius, and the horizon against the episode budget. Failure returns
$\varnothing$ together with a machine-readable reason.

Because nothing is committed that does not verify, the construction can be installed ahead
of an existing solver ladder as a zeroth rung: if it succeeds the ladder is never entered,
and if it fails the ladder runs exactly as it would have. It can add solutions; it cannot
remove one.

\section{Experimental Setup}

\subsection{Robot}
All robots are the \texttt{unicycle2\_v0} model of the dynobench benchmark
\cite{dynobench}, with the published parameters: $v_{\max} = 0.5$\,m/s,
$\omega_{\max} = 0.5$\,rad/s, $a_{\max} = 0.25$\,m/s$^2$,
$\alpha_{\max} = 0.25$\,rad/s$^2$, and a $0.5 \times 0.25$\,m rectangular body
($d_{\mathrm{diag}} = 0.559$\,m). Two consequences of these bounds recur below: the braking
distance is $v_{\max}^2/2a_{\max} = 0.50$\,m, and the minimum turning radius at full speed
is $v_{\max}/\omega_{\max} = 1.00$\,m. Simulation uses $dt = 0.1$\,s. We integrate with
RK4 where dynobench specifies Euler; RK4 is strictly more accurate, so our trajectories are
not bit-comparable with published runs on the same model.

\subsection{Scenarios}
All scenarios use a $17 \times 17$\,m workspace, $T_{\max} = 1300$ steps, goal radius
$0.2$\,m and clearance $\varepsilon = 0.05$\,m. Every robot's goal is another robot's
start, so every instance is a permutation swap.

\begin{itemize}
\item \textbf{open\_cross\_$N$}: $N/2$ facing rows spaced $1.0$\,m apart, each pair
swapping across a $15$\,m traverse. No obstacles.
\item \textbf{cluttered\_cross\_$N$}: the same crossing structure with four box pillars.
Rows are $2.14$\,m apart at $N{=}16$ and $1.0$\,m at $N{=}32$.
\item \textbf{circular\_cross\_$N$}: $N$ robots on a circle of radius $7.5$\,m, each
swapping with its antipode. Every route is a diameter.
\end{itemize}

\subsection{Metrics and caveats}
We report whether the construction succeeded, goals reached and collisions from the
verified rollout, horizon in steps against the budget, the minimum surface-to-surface gap
attained between any two bodies, the device chosen, and the number of solver calls. Three
caveats belong on every number in this paper:

\begin{enumerate}
\item $n = 1$ per scenario. Our instances are deterministic under a fixed initialisation,
so a success here is one sample of one, not a success \emph{rate}.
\item Wall-clock times are not comparable to published figures for the baselines. Ours is
Python driving IPOPT through CasADi; the reference implementations are C++.
\item The baseline is our own reimplementation of the K-ARC ladder, as no code has been
released. Every switch is at the setting the paper describes, but it is our reading of the
paper that is being measured, not the authors' system.
\end{enumerate}

\section{Results}

\subsection{Main result}

Table~\ref{tab:main} gives the outcome on all five scenarios. Four are solved outright with
zero solver calls: no trajectory optimisation, no sampling rung, no conflict-resolution
round. Every solved row is a verified rollout---all goals reached, zero collisions, horizon
inside budget, and a positive surface gap.

\begin{table}[t]
\caption{Constructive planner, all five scenarios. Verified by rollout through the
simulator's own collision checker. ``Gap'' is the minimum surface-to-surface distance
attained between any two bodies.}
\label{tab:main}
\centering
\small
\begin{tabular}{@{}lccccc@{}}
\toprule
Scenario & Solved & Goals & Coll. & Steps/1300 & Gap (m) \\
\midrule
open\_cross\_32      & yes & 32/32 & 0 & 633  & 0.2493 \\
cluttered\_cross\_16 & yes & 16/16 & 0 & 906  & 0.0530 \\
cluttered\_cross\_32 & \textbf{no} & --- & --- & --- & --- \\
circular\_cross\_16  & yes & 16/16 & 0 & 836  & 0.0856 \\
circular\_cross\_32  & yes & 32/32 & 0 & 1212 & 0.0503 \\
\bottomrule
\end{tabular}
\end{table}

\begin{table}[t]
\caption{Device selected per scenario, and the schedule it required.}
\label{tab:device}
\centering
\small
\begin{tabular}{@{}lccc@{}}
\toprule
Scenario & Device & Delayed & Wall (s) \\
\midrule
open\_cross\_32      & 2 lanes, pitch 0.500\,m & 0 (max 0)   & 22 \\
cluttered\_cross\_16 & 2 lanes, pitch 0.345\,m & 6 (max 225) & 8 \\
circular\_cross\_16  & orbit, $\rho=1.551$\,m  & 9 (max 255) & 7 \\
circular\_cross\_32  & orbit, $\rho=3.102$\,m  & 23 (max 475)& 51 \\
\bottomrule
\end{tabular}
\end{table}

The open cross is the cleanest case: 16 independent head-on pairs, two lanes, \emph{zero}
robots delayed, and a $0.2493$\,m surface gap---five times the required clearance---in 633
of 1300 steps. The coordination is entirely spatial; the schedule does nothing.

\subsection{Evidence for device selection}
\label{sec:res-device}

Table~\ref{tab:spread} reports the quantity that drives \eqref{eq:hub}. The separation is
categorical rather than marginal, which is why a single threshold suffices.

\begin{table}[t]
\caption{Spread of pairwise encounter midpoints. The circular scenarios are not ``the open
cross bent into a circle''; every route is a diameter, so all $N$ encounters share one
point.}
\label{tab:spread}
\centering
\small
\begin{tabular}{@{}lcc@{}}
\toprule
Scenario & Midpoint spread (std, m) & Structure \\
\midrule
open\_cross\_32      & $[0.00,\ 4.61]$ & 16 separate crossings \\
cluttered\_cross\_16 & $[0.00,\ 4.91]$ & 8 separate crossings \\
circular\_cross\_16  & $[0.00,\ 0.00]$ & one junction \\
circular\_cross\_32  & $[0.00,\ 0.00]$ & one junction \\
\bottomrule
\end{tabular}
\end{table}

Before the junction device existed, greedy insertion seated \textbf{1 of 32} robots on
circular\_cross\_32. No amount of scheduling could have repaired this: with only two sides
available, thirty-two routes through one point admit no lane assignment at all.

\subsection{Ring radius}

Equation~\eqref{eq:ring} gives the smallest ring on which the component fits. It is also,
empirically, the best (Table~\ref{tab:pack}): inflating it monotonically \emph{reduces} the
number of robots that can be seated. A wider ring is a longer detour through the same
congested annulus, so the extra room costs more time in the contested region than the
separation buys. We found this counter-intuitive and report it because a reader sizing such
a ring would plausibly add margin by default.

\begin{table}[t]
\caption{Ring radius scaling on circular\_cross\_32. Scale $1.0$ is \eqref{eq:ring}.}
\label{tab:pack}
\centering
\small
\begin{tabular}{@{}ccc@{}}
\toprule
Scale & Radius (m) & Robots seated \\
\midrule
\textbf{1.0} & 3.10 & \textbf{32 / 32} \\
1.3 & 4.03 & 21 / 32 \\
1.6 & 4.96 & 18 / 32 \\
2.0 & 6.20 & 18 / 32 \\
\bottomrule
\end{tabular}
\end{table}

\subsection{Scheduling repairs}
\label{sec:res-sched}

Table~\ref{tab:sched} isolates the two repairs on cluttered\_cross\_32, the hardest
instance. Retrying under many random orders saturates well short of a solution, which is
evidence that the difficulty is structural rather than a matter of luck; eviction, which
addresses the measured failure mode directly, does considerably better than 250 random
restarts at a fraction of the cost.

\begin{table}[t]
\caption{Robots seated on cluttered\_cross\_32 by scheduling strategy.}
\label{tab:sched}
\centering
\small
\begin{tabular}{@{}lc@{}}
\toprule
Strategy & Robots seated \\
\midrule
Greedy insertion, 8 candidate orders   & 21 / 32 \\
Greedy insertion, 250 candidate orders & 26 / 32 \\
+ single-blocker eviction              & \textbf{28 / 32} \\
\bottomrule
\end{tabular}
\end{table}

Both repairs had to be held back rather than led with, for the same reason: each rescues
orders that plain greed cannot seat, and the rescued schedule is not necessarily good.
Leading with the precedence order took circular\_cross\_32 from 1212 steps to 1300---its
entire budget---and allowing eviction on the first candidate order did the same. Sweeping
all orders without repairs first, and only then permitting them, recovers the shorter
schedule while keeping the harder instance's gains.

\section{Failure Analysis: cluttered\_cross\_32}
\label{sec:failure}

We report this scenario in detail because the process of diagnosing it produced two wrong
answers before a supported one, and both wrong answers were ours.

\paragraph{Retraction 1: ``not the time budget''} We initially reported that quadrupling
the horizon changed nothing. That measurement was taken on a superseded version of the
scheduler and does not survive re-measurement: horizon does buy improvement before
saturating.

\paragraph{Retraction 2: ``robots park on each other's routes''} We initially attributed
the failure to robots resting on paths others needed. This came from a diagnostic of our own
that classified a blocker as ``parked'' only if it blocked at \emph{every} candidate delay,
which mislabels the common case. Sharpened to report where the blocker actually stands, the
blockers are \emph{driving}. The supporting geometry agrees: of 40 route/goal incursions in
this scenario, \textbf{zero} are mid-route---every one is at an endpoint, which is
unavoidable in a permutation swap and equally true of cluttered\_cross\_16, which solves.

\paragraph{The supported diagnosis} cluttered\_cross\_32 combines open\_cross\_32's row
density ($1.0$\,m) with four pillars; cluttered\_cross\_16 has the same pillars but
$2.14$\,m rows. The pillars force detours, the detours cross rows, and the conflict graph
grows from 16 pairs to 52 (mean degree $3.25$) with a single 20-robot component. That
component's colouring demands four lanes, and the corridor holds three
(Table~\ref{tab:capacity}). The residue is a genuine capacity shortfall: more independent
encounters than the corridor has lanes, and too little horizon to serialise the remainder.

\begin{table}[t]
\caption{Lane demand against corridor capacity.}
\label{tab:capacity}
\centering
\small
\begin{tabular}{@{}lccc@{}}
\toprule
Scenario & Lanes wanted & Span needed & Corridor \\
\midrule
open\_cross\_32      & 2 & 0.50\,m & 0.70\,m \\
cluttered\_cross\_32 & 4 & \textbf{1.50\,m} & \textbf{0.70\,m} \\
\bottomrule
\end{tabular}
\end{table}

That the same 20-robot component is also the only component in any scenario whose route
directions disagree (tangent eigenvalue ratio $0.159$, against $\le 0.038$ everywhere else)
motivated the direction-grouping of Sec.~\ref{sec:lane}, which is correct in itself but did
not change the outcome here.

\paragraph{A device that did not transfer} Every pillar in this scenario is passed on both
sides ($4/8$, $5/10$, $4/10$, $5/8$ robots), and each such split turns the narrow gap beside
a pillar into a head-on encounter---apparently the same failure the junction device
addresses. Applying the right-hand rule to pillars made matters substantially worse.
Normalising every route onto a ring around each pillar took seating from 28 to \textbf{4}
and stopped cluttered\_cross\_16 solving at all; flipping only the wrong-side minority still
gave \textbf{12}. A forced flip around a $4$\,m pillar is a long detour that pushes the
route into neighbouring rows, costing more than the encounters it removes. The lesson worth
carrying: the right-hand rule works when routes genuinely converge on a point, and it does
not follow that it works for an obstacle they merely pass.

\section{Limitations}
\label{sec:lim}

\subsection{Progress discretisation}
As set out in Sec.~\ref{sec:refs}, conflicts are detected on a 64-sample arclength grid,
comparing robots at equal progress rather than equal time. The realised trajectories are
continuous-time integrations of \eqref{eq:dyn} and the acceptance test is a full rollout, so
this approximation cannot produce an unsafe accepted plan---an unsound conflict set yields a
bad \emph{assignment}, which verification then rejects. But it does bound the method's
reach: where robots progress at very different rates, the proxy degrades and the
construction will be rejected more often than it need be. Replacing the grid with a
continuous conflict predicate over the realised trajectories is the most valuable single
extension we see, and would remove the one genuinely discrete step in an otherwise
continuous pipeline.

\subsection{Interleaved control}
The rest-to-rest decomposition of Sec.~\ref{sec:legs} never fires linear and angular
acceleration together: across all three solved scenarios, \textbf{zero} of the
$10\,057$--$25\,808$ control steps have both $a \ne 0$ and $\alpha \ne 0$. Nothing in
\eqref{eq:dyn} forbids it; the restriction buys closed-form feasibility and the guarantee
that the executed path is the polyline the clearance reasoning used. Its cost is real
(Table~\ref{tab:control}): 21--33\% of motion steps are spent turning in place, robots come
to a full stop 3--6 times mid-route, and the horizon is $1.70$--$2.08\times$ a
cruise-limited lower bound for the same path length.

\begin{table}[t]
\caption{Cost of the rest-to-rest decomposition. ``Both'' counts control steps with $a$ and
$\alpha$ simultaneously non-zero.}
\label{tab:control}
\centering
\small
\begin{tabular}{@{}lcccc@{}}
\toprule
Scenario & Both & Turning & Stops & vs.\ bound \\
\midrule
open\_cross\_32      & 0 / 17\,324 & 21.2\% & 3 & 1.70$\times$ \\
cluttered\_cross\_16 & 0 / 10\,057 & 26.0\% & 4 & 1.84$\times$ \\
circular\_cross\_32  & 0 / 25\,808 & 33.4\% & 6 & 2.08$\times$ \\
\bottomrule
\end{tabular}
\end{table}

The obvious remedy---arc through corners with both controls active---is constrained by the
robot rather than by the construction. Cornering on radius $\rho$ requires
$v \le \omega_{\max}\rho$, so at $v_{\max}$ the minimum radius is $1.00$\,m while lanes are
$0.35$--$0.50$\,m wide. A robot cannot corner at speed inside its own lane. There is
nonetheless clear headroom at the small end: on open\_cross\_32 the median turn is
$6.7^\circ$ and 62\% of turns are under $30^\circ$, artefacts of route simplification for
which a full stop is plainly wasteful. Steering through small turns while driving, and
retaining the stop-and-go decomposition for genuine corners, would recover much of the
$1.70\times$ without disturbing the clearance argument, since the lateral deviation from
the chord for a turn $\delta$ over a leg of length $L$ is approximately $L\delta/8$---about
$0.05$\,m for $7^\circ$ over $5$\,m, well inside a lane.

\subsection{Incompleteness and scope}
Insertion scheduling is greedy and offers no completeness guarantee; the repairs of
Sec.~\ref{sec:sched} are bounded and targeted, not a search. Device selection currently
distinguishes two cases, and cluttered\_cross\_32 suggests at least one more is needed for
components that exceed corridor capacity. All results are $n = 1$ per scenario, on one robot
model, in one workspace size.

\section{Conclusion}

We have argued that dense kinodynamic multi-robot planning is better served by deciding the
coordination first than by repeatedly asking a planner for trajectories, and that the
conflict graph's own geometry is sufficient to choose the coordination. A component of
scattered crossings takes lanes; a component whose encounters converge on a point takes a
roundabout; and the distinction is measurable, categorical in our benchmarks, and decides
which of two structurally incompatible displacement rules applies. Four of five dense
scenarios with 16--32 second-order unicycles are solved with zero solver calls in seconds
of wall time, each verified by rollout, and installed ahead of an existing ladder the
construction can only add solutions.

The most valuable next steps follow directly from the limitations. Replacing the progress
grid with a continuous conflict predicate removes the one discrete step in the pipeline.
Steering through small turns removes most of a measured $1.70$--$2.08\times$ penalty. And
cluttered\_cross\_32 is a standing request for a third device: one that adds coordination
capacity without lateral room, for components whose lane demand exceeds what the corridor
can hold.

\section*{Disclosure}

An AI assistant (Claude, Anthropic) was used during this work for implementation,
experiment execution, and drafting. It is not an author. All experimental claims in this
paper are reproducible from the accompanying code and configuration; all references were
checked against their primary sources before inclusion. The retractions in
Sec.~\ref{sec:failure} correct claims made earlier in the same development process.

\bibliographystyle{IEEEtran}
\bibliography{refs}

\end{document}
